\subsection{Vzdálenost bodu od přímky}
\label{subsec:Distance_point_line_space}

Vzdálenost dvou bodů v prostoru jsme již určovali pomocí délky
vektoru (viz podkapitola \ref{subsubsec:Distance_between_points_in_space}) který tyto body spojuje. Nyní budeme hledat vzdálenost bodu
od přímky.

Je-li dán bod $P$ a přímka $p$, můžeme bod $P$ spojit s libovolným
bodem přímky $p$. Takových úseček existuje nekonečně mnoho, jejich
délky však obecně nejsou stejné. Vzdáleností bodu od přímky rozumíme
délku nejkratší z nich.

Nejkratší spojnice bodu $P$ s přímkou $p$ je na přímku $p$ kolmá.

\begin{figure}[H]
    \centering
    \begin{tikzpicture}[scale=0.9]

        % Body
        \coordinate (A) at (0.5,0);
        \coordinate (U) at (3.5,0);
        \coordinate (H) at (4.5,0);
        \coordinate (P) at (4.5,2.8);
        \coordinate (PU) at (7.5,2.8);

        % Rovnoběžník
        \fill[gray!10]
            (A) -- (U) -- (PU) -- (P) -- cycle;

        \draw[gray!60]
            (U) -- (PU) -- (P);

        % Přímka
        \draw[thick]
            (-0.5,0) -- (8.2,0)
            node[right] {$p$};

        % Vektory
        \draw[vector]
            (A) -- (U)
            node[midway, below] {$\vec{u}$};

        \draw[vector]
            (A) -- (P)
            node[midway, above left] {$\overrightarrow{AP}$};

        % Kolmice
        \draw[dashed, thick]
            (P) -- (H)
            node[midway, right] {$d(P,p)$};

        % Pravý úhel
        \draw
            (H) ++(-0.22,0)
            -- ++(0,0.22)
            -- ++(0.22,0);

        % Body
        \fill (A) circle (2pt) node[below] {$A$};
        \fill (H) circle (2pt) node[below] {$H$};
        \fill (P) circle (2pt) node[above] {$P$};

        \node at (6.1,1.4) {$S$};

    \end{tikzpicture}
    \caption{Vzdálenost bodu $P$ od přímky $p$.}
    \label{fig:Distance_point_line_space}
\end{figure}

\begin{definition}
    \label{def:Distance_point_line_space}
    \textbf{Vzdáleností bodu $P$ od přímky $p$} nazýváme nejmenší
    vzdálenost bodu $P$ od libovolného bodu přímky $p$.

    Označíme-li $H$ patu kolmice vedené z bodu $P$ na přímku $p$,
    potom
    \[
        d(P,p)=|PH|.
    \]

    Ekvivalentně můžeme psát
    \[
        d(P,p)
        =
        \min_{X\in p}|PX|.
    \]
\end{definition}

Nyní odvodíme vztah, pomocí kterého můžeme tuto vzdálenost určit
přímo ze souřadnic.

Uvažujme přímku $p$, která prochází bodem
\[
    A[x_A;y_A;z_A]
\]
a má nenulový směrový vektor
\[
    \vec{u}=(u_x;u_y;u_z).
\]

Dále mějme bod
\[
    P[x_P;y_P;z_P].
\]

Z bodu $A$ do bodu $P$ vede vektor
\[
    \overrightarrow{AP}
    =
    (x_P-x_A;y_P-y_A;z_P-z_A).
\]

Vektory $\overrightarrow{AP}$ a $\vec{u}$ určují rovnoběžník.
Jeho obsah můžeme podle geometrického významu vektorového součinu (viz podkapitola \ref{subsubsec:Vector_cross_product_in_space})
vyjádřit jako
\[
    S
    =
    \left|
        \overrightarrow{AP}\times\vec{u}
    \right|.
\]

Na druhou stranu můžeme obsah stejného rovnoběžníku vypočítat jako
součin délky jeho základny a příslušné výšky. Za základnu zvolíme
vektor $\vec{u}$. Výškou rovnoběžníku je právě vzdálenost bodu $P$
od přímky $p$, neboli $|PH|=d(P,p)$.

Platí tedy
\[
    S
    =
    |\vec{u}|\cdot d(P,p).
\]

Porovnáním obou vztahů dostáváme
\[
    |\vec{u}|\cdot d(P,p)
    =
    \left|
        \overrightarrow{AP}\times\vec{u}
    \right|,
\]
a tedy

\begin{equation}
        d(P,p)
        =
        \frac{
            \left|
                \overrightarrow{AP}\times\vec{u}
            \right|
        }{
            |\vec{u}|
        }
    \label{eq:Distance_point_line_space}
\end{equation}

K výpočtu vzdálenosti bodu od přímky tedy potřebujeme libovolný bod
$A$ ležící na přímce $p$, její směrový vektor $\vec{u}$ a vektor
$\overrightarrow{AP}$ vedoucí z bodu přímky do zadaného bodu.

\begin{remark}
    \label{rem:Choice_point_direction_distance_line}
    Ve vztahu
    \[
        d(P,p)
        =
        \frac{
            |\overrightarrow{AP}\times\vec{u}|
        }{
            |\vec{u}|
        }
    \]
    můžeme zvolit \textbf{libovolný bod $A$ přímky $p$} a libovolný
    její nenulový směrový vektor.

    Změníme-li směrový vektor například na $k\vec{u}$, kde $k\neq0$,
    změní se čitatel i jmenovatel ve stejném poměru:
    \[
        \frac{
            |\overrightarrow{AP}\times(k\vec{u})|
        }{
            |k\vec{u}|
        }
        =
        \frac{
            |k|\,
            |\overrightarrow{AP}\times\vec{u}|
        }{
            |k|\,|\vec{u}|
        }.
    \]
    Výsledná vzdálenost se proto nezmění.
\end{remark}

\begin{example}
    \label{exa:Distance_point_parametric_line_space}

    Určete vzdálenost bodu $P[2;2;4]$ od přímky
    \[
        \begin{aligned}
            p:  x&=1+2t,\\
                y&=-1+t,\\
                z&=2+2t; t\in\mathbb{R}.
        \end{aligned}
    \]

    Z parametrického vyjádření přímky určíme bod
    \[
        A[1;-1;2]
    \]
    a její směrový vektor
    \[
        \vec{u}=(2;1;2).
    \]

    Vektor z bodu $A$ do bodu $P$ je
    \[
        \overrightarrow{AP}
        =
        P-A
        =
        (1;3;2).
    \]

    Vypočítáme vektorový součin
    \begin{align*}
        \overrightarrow{AP}\times\vec{u}
        &=
        (1;3;2)\times(2;1;2)\\
        &=
        \left(
            3\cdot2-2\cdot1;
            2\cdot2-1\cdot2;
            1\cdot1-3\cdot2
        \right)\\
        &=
        (4;2;-5).
    \end{align*}

    Jeho velikost je
    \[
        \left|
            \overrightarrow{AP}\times\vec{u}
        \right|
        =
        \sqrt{4^2+2^2+(-5)^2}
        =
        \sqrt{45}
        =
        3\sqrt5.
    \]

    Velikost směrového vektoru je
    \[
        |\vec{u}|
        =
        \sqrt{2^2+1^2+2^2}
        =
        3.
    \]

    Dosadíme do vztahu pro vzdálenost bodu od přímky:
    \[
        d(P,p)
        =
        \frac{
            |\overrightarrow{AP}\times\vec{u}|
        }{
            |\vec{u}|
        }
        =
        \frac{3\sqrt5}{3}
        =
        \sqrt5.
    \]

    Vzdálenost bodu $P$ od přímky $p$ je tedy $d(P,p)=\sqrt{5}$.
\end{example}

Vzdálenost bodu od přímky můžeme určit také jiným způsobem.
Místo výpočtu obsahu rovnoběžníku můžeme přímo najít bod přímky,
který je k bodu $P$ nejblíže.

Označme tento bod $H$. Protože $H$ leží na přímce $p$, můžeme jej
zapsat ve tvaru
\[
    H=A+t\vec{u}.
\]

Zároveň musí být úsečka $PH$ kolmá k přímce $p$. Vektor
$\overrightarrow{PH}$ je tedy kolmý ke směrovému vektoru $\vec{u}$:
\[
    \overrightarrow{PH}\cdot\vec{u}=0.
\]

Protože
\[
    \overrightarrow{PH}
    =
    H-P
    =
    A+t\vec{u}-P,
\]
dostáváme
\[
    (A+t\vec{u}-P)\cdot\vec{u}=0.
\]

Po roznásobení
\[
    (A-P)\cdot\vec{u}
    +
    t|\vec{u}|^2
    =
    0,
\]
a tedy
\[
    t
    =
    \frac{
        \overrightarrow{AP}\cdot\vec{u}
    }{
        |\vec{u}|^2
    }.
\]

Tato hodnota parametru určuje \textbf{patu kolmice} $H$, tedy bod
přímky $p$, který leží k bodu $P$ nejblíže. Vzdálenost potom můžeme
určit jako délku úsečky $PH$.

\begin{example}
    \label{exa:Foot_perpendicular_point_line_space}

    Vraťme se k předchozímu příkladu. Přímka
    \[
        \begin{aligned}
            p: x&=1+2t,\\
            y&=-1+t,\\
            z&=2+2t; t\in\mathbb{R}
        \end{aligned}
    \]
    má směrový vektor
    \[
        \vec{u}=(2;1;2)
    \]
    a hledáme bod $H\in p$ takový, aby
    \[
        PH\perp p.
    \]

    Obecný bod přímky je
    \[
        H[1+2t;-1+t;2+2t].
    \]

    Pro bod $P[2;2;4]$ dostáváme vektor
    \[
        \overrightarrow{PH}
        =
        (-1+2t;-3+t;-2+2t).
    \]

    Protože $\overrightarrow{PH}\perp\vec{u}$, musí platit
    \[
        \overrightarrow{PH}\cdot\vec{u}=0.
    \]

    Dosadíme:
    \begin{align*}
        (-1+2t)\cdot2
        +(-3+t)\cdot1
        +(-2+2t)\cdot2
        &=0,\\
        -2+4t-3+t-4+4t
        &=0,\\
        9t-9&=0,\\
        t&=1.
    \end{align*}

    Pata kolmice má tedy souřadnice
    \[
        H[3;0;4].
    \]

    Vzdálenost bodu $P$ od přímky $p$ je
    \[
        d(P,p)
        =
        |PH|
        =
        \sqrt{(3-2)^2+(0-2)^2+(4-4)^2}
        =
        \sqrt5.
    \]

    Dostáváme stejný výsledek jako při použití vektorového součinu:
    \[
        d(P,p)=\sqrt5.
    \]
\end{example}

\begin{remark}
    \label{rem:Point_on_line_distance_zero_space}
    Leží-li bod $P$ přímo na přímce $p$, je jeho vzdálenost od přímky
    rovna nule.

    Také vztah
    \[
        d(P,p)
        =
        \frac{
            |\overrightarrow{AP}\times\vec{u}|
        }{
            |\vec{u}|
        }
    \]
    tuto skutečnost zachycuje. Pokud $P\in p$, jsou vektory
    $\overrightarrow{AP}$ a $\vec{u}$ rovnoběžné, takže
    \[
        \overrightarrow{AP}\times\vec{u}
        =
        \vec{0},
    \]
    a proto
    \[
        d(P,p)=0.
    \]
\end{remark}

\begin{example}
    \label{exa:Distance_point_line_parameter_space}

    Je dána soustava přímek
    \[
        \begin{aligned}
            p_m: x&=t,\\
            y&=m,\\
            z&=t; t\in\mathbb{R}.
        \end{aligned}
    \]

    Určete hodnoty parametru $m\in\mathbb{R}$ tak, aby vzdálenost
    přímky $p_m$ od počátku soustavy souřadnic $O[0;0;0]$ byla rovna $3$.

    Bodem přímky může být
    \[
        A[0;m;0]
    \]
    a jejím směrovým vektorem je
    \[
        \vec{u}=(1;0;1).
    \]

    Vektor
    \[
        \overrightarrow{AO}
        =
        O-A
        =
        (0;-m;0).
    \]

    Jejich vektorový součin je
    \[
        \overrightarrow{AO}\times\vec{u}
        =
        (0;-m;0)\times(1;0;1)
        =
        (-m;0;m).
    \]

    Proto
    \[
        \left|
            \overrightarrow{AO}\times\vec{u}
        \right|
        =
        \sqrt{m^2+m^2}
        =
        |m|\sqrt2
    \]
    a
    \[
        |\vec{u}|=\sqrt2.
    \]

    Vzdálenost počátku od přímky je tedy
    \[
        d(O,p_m)
        =
        \frac{|m|\sqrt2}{\sqrt2}
        =
        |m|.
    \]

    Požadujeme
    \[
        |m|=3,
    \]
    odkud
    \[
        m=\pm3,
    \]
    tedy přímky $p_3$ a $p_{-3}$ mají od počátku vzdálenost $3$ a jejich předpisy jsou:
    \[
        \begin{aligned}
            p_3: x&=t,\\
                 y&=3,\\
                 z&=t; t\in\mathbb{R}
        \end{aligned}
        \qquad
        \begin{aligned}
            p_{-3}: x&=t,\\
                    y&=-3,\\
                    z&=t; t\in\mathbb{R}.
        \end{aligned}
    \]
\end{example}